\chapter{Calculus and Analysis For Machine Learning}

Possible sections:
\begin{itemize}
  \item Some basic calculus results (MVT, extreme value theorem, etc.) 
  \item Fundamental theorem of calculus (then, implications? Like, how is this useful for...)
  \item Inverse and implicit function theorems (and implications) 
  \item Divergence theorem (and relationship to fundamental theorem of calculus)
  \item Stokes theorem? (Does this apply to higher than $\R^3$?) 
  \item integration by parts (and some implications. Like for ML?) 
  \item Multivariable chain rule
  \item Change of variables theorem
\end{itemize}


\subsection{The Laplacian}
The Laplacian is foundational to calculus in higher dimensions. This section attempts to build some intuition about it. In particular, how it can be viewed as a local averaging operator. Let $f\in C^2$. We claim that the following holds:
\begin{equation} \label{eq:laplace-local-average}
\frac{1}{n} \Delta f(x) = \frac{1}{r^2}\left(\int_{S_n} f(z) d\mu(z) - f(x) \right) + O(r^2). 
\end{equation}
We're going to build some intuition for what this formula means and prove that it holds. 

To build intuition, we start with the 1d case. In 1 dimension, claim that
\begin{equation} \label{eq:2nd-deriv-avg-1d}
f''(x) = \frac{f(x+r) + f(x-r) - 2f(x) }{2r^2} + O(r^2). 
\end{equation}
To see that this holds, consider the Taylor expansion about $f$
$$
f(x+r) = f(x) + rf'(x) + \frac{1}{2}r^2 f'(x) + \frac{1}{3}r^3 f'''(x) + O(r^4).
$$
Note that the Taylor series of f(x-r) computed to the same order is identical except for the fact that odd terms have a negative sign. Adding the Taylor series for $f(x+r)$ and $f(x-r)$ we see that
$$
f(x+r) + f(x-r) = 2f(x) + r^2f''(x) + O(r^4). 
$$
Rearranging terms and dividing both sides by $r^2$ we get \eqref{eq:2nd-deriv-avg-1d}. 

Some observations:
This equation may seem intuitively obvious from the definition of the derivative, e.g., the definition of the Frechet derivative. It seems more obvious to me that way. Apply the definition of Frechet derivative twice. Instead of a tangent line, you're fitting a tangent quadratic. However, once you start thinking like that, you're thinking in terms of Taylor series. The n-th derivative is really just the best n'th order polynomial fit. Well, nearly. You have to normalize for 1/n!. But that's kind of it.

What confuses me here is how to think of the role of r as you take the limit. In the definition of a derivative, I get it. I think. You take a line passing through two points. And then bring the points together and look at the limit. Here, I guess it's the same? You are fitting a quadratic (this is the coefficient of the quadratic?) by taking the value at three points: x and two neighboring points. Then taking the limit. I'm confused about how to think about the local average operation though. Like, when I think of the behavior of the Laplacian operating in a PDE, I guess I can kind of think of it like, at each point, it sort of looks in a neighborhood. And checks its value relative to neighbors. And the $1/r^2$ is the scaling factor that makes that intuition work in small neighborhoods as you take a limit. 

...yep, I don't get this. Like, if you take one derivative you divide by r. But, up top, it doesn't seem that different from one derivative. You're taking the difference $f(x+r) - f(x)$. But also, I guess $f(x-r) - f(x)$ which is like a linear measurement in the backwards direction. And then you blow them both up by dividing by $r^2$? If you divide by $r^2$ in a traditional derivative, then you really blow the whole thing up. I mean, assuming the derivative exists, then it has to be finite, so dividing by $r^2$ has to blow it up. Why doesn't it here? I'm just kind of confused. 

OK, regardless, the local averaging intuition seems plausible. It's somehow like it's looking in a small neighborhood and then taking a local average. 

Moving on to the general case in $n$ dimensions. Based on the intuition above, it seems natural to prove this using a Taylor series. That's probably a good exercise. I'm going to put a proof down here that's slicker that uses the divergence theorem. And I guess Taylor's theorem at the end. 

The equation \eqref{eq:laplace-local-average} shows how this intuition extends to higher dimensions. The main difference between \eqref{eq:laplace-local-average}  and \eqref{eq:2nd-deriv-avg-1d}. is the $n$ term on the left hand side. Other than this, the expressions are fundamentally identical. The Laplacian is just the local average over $r^2$ + $O(r^2)$. The $n$ term indicates what happens when we add more dimensions. In the Laplacian we're taking the sum second partial derivatives. The $n$ term normalizes this so the averaging relationship still holds. One may think of it as having to consider the mean local curvature rather than the sum of local curvatures. 

\begin{proof}
We now prove the general case for $n$ dimensions \eqref{eq:laplace-local-average}. This proof is from stack exchange. I think it's really slick. I haven't fully wrapped my head around it. 

Given a point $p$ in the domain of the function $f$, consider the function which is the average value $A(r)$ of $f$ over the sphere $S_r$ of radius $r$ centered at $p$. If $S^{n-1}$ is the unit sphere in $\mathbb{R}^n$ then you can write the function as:

$$A(r) = \int_{S^{n-1}} f(p + rx) \, dx$$

I'm integrating with respect to the standard content on the sphere, but rescaled so that the sphere has unit content. Using Leibnitz integral rule, differentiating $A(r)$ with respect to $r$ gives

$$\frac{dA}{dr} = \int_{S^{n-1}} x \cdot \nabla f(p + rx) \, dx.$$

In the above equation $x \in S^{n-1}$ is a unit vector, and the $\cdot$ is dot product of vectors. The above can be interpreted as a flux integral, so Gauss's theorem about flux integrals applies. Before applying it, first note that\footnote{THIS IS MESSED UP. And some stuff below is messed up. Like, we're integrating over the unit ball below. Which seems right. But I don't know that we're handling the divergence theorem properly before that. Because you're thinking about integrating over a scaled volume in the divergence theorem. Maybe I should think of it as scaling the volume, applying the divergence theorem, then going back to the unit sphere for convenience when we take a limit in r?}
$$
\nabla \cdot \nabla f(p+rx) = r\Delta f(p+rx).
$$
This follows from the definition of the Laplacian as a sum of partial derivatives and the chain rule. Now, applying the divergence theorem we have

$$\frac{dA}{dr} = r \int_{D^n} \nabla^2 f(p + rx) \, dx$$

where $x$ is now a variable for the unit ball $D^n$. So this tells you the first few terms of the Taylor expansion for the function $A$, in particular $A(0) = f(p)$, $A'(0) = 0$, and so on with the next terms being some multiple of $\nabla^2 f(p)$. A quick calculation says it should be $\frac{\nabla^2 f(p)}{n}$. TODO: Explain each of these last terms. For the last one, you just take a derivative of the da over dr term and take the limit as r to 0. 
\end{proof}
